\documentclass[conference]{IEEEtran}
\usepackage{amsmath,amssymb,graphicx,booktabs,cite}
\begin{document}
\title{Covariance-Gated Predictive Fire Control for a Laser Turret Engaging a Maneuvering Aerial Target}
\author{\IEEEauthorblockN{Venky (and team)}\IEEEauthorblockA{Dept. of Mechanical Engineering, IIT Bombay\\ Code: \texttt{https://github.com/<user>/me319-guardian}}}
\maketitle

\begin{abstract}
% 120-150 words: problem, EKF-CA estimator, tp=0.5 s prediction, notch-PI-lead Gc(s), gated firing, key results (hit rate, intercept time, Monte-Carlo).
\end{abstract}

\section{Introduction}
% Problem, objectives, assumptions (planar, balanced trunnion, no voltage limit, negligible beam time-of-flight), contribution list.

\section{System Architecture and Design}
% Fig: engineering schematic (turret, radar, LOS, predicted point).
% Fig: closed-loop block diagram: radar -> EKF -> predictor -> theta_ref -> Gc -> Gp; feedforward path; fire gate.
% Interpretation of tp (aim-point vs tracking) and justification.

\section{Mathematical Modeling}
\subsection{Turret Elevation Plant}
Reflecting the rotor inertia through the $N{:}1$ gearbox, $J_{eq}=J_t+N^2J_m=3200$~kg\,m$^2$,
\begin{align}
J_{eq}\ddot\theta + b\dot\theta &= N K_t i, \\
L_m \dot i + R_m i &= V_m - N K_e \dot\theta,
\end{align}
\begin{equation}
G_p(s)=\frac{\Theta(s)}{V_m(s)}=\frac{NK_t}{s\left[(J_{eq}s+b)(L_ms+R_m)+N^2K_tK_e\right]}
=\frac{3.125}{s\,(s^2+10.47s+3129.7)}.
\end{equation}
% poles 0, -5.23 +/- 55.70j (wn=55.9 rad/s, zeta=0.094); interpretation.
\subsection{Target Estimation Model}
% CA state x=[p_x p_y v_x v_y a_x a_y]^T, white-jerk process noise q, polar measurement h(x), EKF.
\subsection{Prediction and Elevation Reference}
% x(t+tp)=F(tp)x, theta_ref = atan2, LOS rate/accel expressions.
\subsection{Closed-Loop Model}

\section{Controller Design and Analysis}
% P-only gain ceiling; notch + PI + lead; Bode, margins, root locus; Ka; noise -> turret jitter; blackout covariance growth; feedforward.

\section{Simulation and Results}
% Table: scenario A/B/C -- first hit time, hits/shots, mean miss, rms perpendicular error, est. rms.
% Figs: target & estimate, theta vs theta_ref, perpendicular miss vs time with shots, blackout zoom.
% Monte-Carlo R1..R100 summary.

\section{Final Evaluation}
% Instructor cases: table + interpretation.

\section{Conclusions}

\bibliographystyle{IEEEtran}
% \bibliography{refs}

\clearpage
\appendices
\section{Use of AI and Computational Tools}
% Tools: Claude (Anthropic) for ..., MATLAB/Octave for ...
% Accepted / modified / rejected outputs (log as you go!).
% Validation reasoning.
\end{document}
